\documentclass[UTF8,12pt]{ctexart}
\usepackage[a4paper,margin=24mm]{geometry}
\usepackage{amsmath,amssymb,bm,graphicx,adjustbox}
\newcommand{\fitmath}[1]{\begin{adjustbox}{max width=\linewidth}$\displaystyle #1$\end{adjustbox}}
\setlength{\parindent}{0pt}
\setlength{\parskip}{.7em}
\sloppy
\allowdisplaybreaks
\begin{document}
\section*{1996 年考研数学三 · 真题}
\subsection*{1996 年全国硕士研究生招生考试试题（试卷 IV）}

\subsection*{一、填空题（本题共 5 小题，每小题 3 分，满分 15 分）}

（1）设方程 \(\displaystyle x=y^y\) 确定 \(\displaystyle y\) 是 \(\displaystyle x\) 的函数，则 \(\displaystyle dy=\)\_\_\_\_\_\_。


（2）设 \(\displaystyle \int xf(x)\,dx=\arcsin x+C\)，则 \(\displaystyle \int\dfrac{\displaystyle 1}{\displaystyle f(x)}\,dx=\)\_\_\_\_\_\_。


（3）设 \(\displaystyle (x_0,\allowbreak y_0)\) 是抛物线 \(\displaystyle y=ax^2+bx+c\) 上的一点。若在该点的切线过原点，则系数应满足的关系是\_\_\_\_\_\_。


（4）设 \(\displaystyle A=\begin{pmatrix}1&1&1&\cdots&1\\a_1&a_2&a_3&\cdots&a_n\\a_1^2&a_2^2&a_3^2&\cdots&a_n^2\\\vdots&\vdots&\vdots&&\vdots\\a_1^{n-1}&a_2^{n-1}&a_3^{n-1}&\cdots&a_n^{n-1}\end{pmatrix},\allowbreak \ X=(x_1,\allowbreak x_2,\allowbreak \ldots,\allowbreak x_n)^{\mathrm T},\allowbreak \ B=(1,\allowbreak 1,\allowbreak \ldots,\allowbreak 1)^{\mathrm T}\)，其中 \(\displaystyle a_i\ne a_j\)（\(\displaystyle i\ne j;i,\allowbreak j=1,\allowbreak 2,\allowbreak \ldots,\allowbreak n\)），则线性方程组 \(\displaystyle A^{\mathrm T}X=B\) 的解是\_\_\_\_\_\_。


（5）（超纲题）设由来自正态总体 \(\displaystyle X\sim N(\mu,\allowbreak 0.9^2)\) 容量为 9 的简单随机样本，得样本均值 \(\displaystyle \bar X=5\)，则未知参数 \(\displaystyle \mu\) 的置信度为 0.95 的置信区间是\_\_\_\_\_\_。


\subsection*{二、选择题（本题共 5 小题，每小题 3 分，满分 15 分）}

（1）累次积分 \(\displaystyle \int_0^{\pi/2}d\theta\displaystyle \int_0^{\cos\theta}f(r\cos\theta,\allowbreak r\sin\theta)r\,dr\) 可以写成（　）。


（A）\(\displaystyle \int_0^1dy\displaystyle \int_0^{\sqrt{y-y^2}}f(x,\allowbreak y)\,dx\)。 （B）\(\displaystyle \int_0^1dy\displaystyle \int_0^{\sqrt{1-y^2}}f(x,\allowbreak y)\,dx\)。


（C）\(\displaystyle \int_0^1dx\displaystyle \int_0^1f(x,\allowbreak y)\,dy\)。 （D）\(\displaystyle \int_0^1dx\displaystyle \int_0^{\sqrt{x-x^2}}f(x,\allowbreak y)\,dy\)。


（2）下述各选项正确的是（　）。


（A）若 \(\displaystyle \sum_{n=1}^{\infty}u_n^2\) 和 \(\displaystyle \sum_{n=1}^{\infty}v_n^2\) 都收敛，则 \(\displaystyle \sum_{n=1}^{\infty}(u_n+v_n)^2\) 收敛。


（B）若 \(\displaystyle \sum_{n=1}^{\infty}|u_nv_n|\) 收敛，则 \(\displaystyle \sum_{n=1}^{\infty}u_n^2\) 与 \(\displaystyle \sum_{n=1}^{\infty}v_n^2\) 都收敛。


（C）若正项级数 \(\displaystyle \sum_{n=1}^{\infty}u_n\) 发散，则 \(\displaystyle u_n\geq\dfrac{\displaystyle 1}{\displaystyle n}\)。


（D）若级数 \(\displaystyle \sum_{n=1}^{\infty}u_n\) 收敛，且 \(\displaystyle u_n\geq v_n\)（\(\displaystyle n=1,\allowbreak 2,\allowbreak \ldots\)），则级数 \(\displaystyle \sum_{n=1}^{\infty}v_n\) 也收敛。


（3）设 \(\displaystyle n\) 阶矩阵 \(\displaystyle A\) 非奇异（\(\displaystyle n\geq2\)），\(\displaystyle A^*\) 是矩阵 \(\displaystyle A\) 的伴随矩阵，则（　）。


（A）\(\displaystyle (A^*)^*=|A|^{n-1}A\)。 （B）\(\displaystyle (A^*)^*=|A|^{n+1}A\)。 （C）\(\displaystyle (A^*)^*=|A|^{n-2}A\)。 （D）\(\displaystyle (A^*)^*=|A|^{n+2}A\)。


\subsection*{试卷 IV（续）}

二、（4）设有任意两个 \(\displaystyle n\) 维向量组 \(\displaystyle \alpha_1,\allowbreak \ldots,\allowbreak \alpha_m\) 和 \(\displaystyle \beta_1,\allowbreak \ldots,\allowbreak \beta_m\)，若存在两组不全为零的数 \(\displaystyle \lambda_1,\allowbreak \ldots,\allowbreak \lambda_m\) 和 \(\displaystyle k_1,\allowbreak \ldots,\allowbreak k_m\)，使 \(\displaystyle (\lambda_1+k_1)\alpha_1+\cdots+(\lambda_m+k_m)\alpha_m+(\lambda_1-k_1)\beta_1+\cdots+(\lambda_m-k_m)\beta_m=0\)，则（　）。


（A）\(\displaystyle \alpha_1,\allowbreak \ldots,\allowbreak \alpha_m\) 和 \(\displaystyle \beta_1,\allowbreak \ldots,\allowbreak \beta_m\) 都线性相关。 （B）\(\displaystyle \alpha_1,\allowbreak \ldots,\allowbreak \alpha_m\) 和 \(\displaystyle \beta_1,\allowbreak \ldots,\allowbreak \beta_m\) 都线性无关。


（C）\(\displaystyle \alpha_1+\beta_1,\allowbreak \ldots,\allowbreak \alpha_m+\beta_m,\allowbreak \alpha_1-\beta_1,\allowbreak \ldots,\allowbreak \alpha_m-\beta_m\) 线性无关。 （D）\(\displaystyle \alpha_1+\beta_1,\allowbreak \ldots,\allowbreak \alpha_m+\beta_m,\allowbreak \alpha_1-\beta_1,\allowbreak \ldots,\allowbreak \alpha_m-\beta_m\) 线性相关。


（5）已知 \(\displaystyle 0<P(B)<1\) 且 \(\displaystyle P[(A_1+A_2)\mid B]=P(A_1\mid B)+P(A_2\mid B)\)，则下列选项成立的是（　）。


（A）\(\displaystyle P[(A_1+A_2)\mid\bar B]=P(A_1\mid\bar B)+P(A_2\mid\bar B)\)。 （B）\(\displaystyle P(A_1B+A_2B)=P(A_1B)+P(A_2B)\)。


（C）\(\displaystyle P(A_1+A_2)=P(A_1\mid B)+P(A_2\mid B)\)。 （D）\(\displaystyle P(B)=P(A_1)P(B\mid A_1)+P(A_2)P(B\mid A_2)\)。


三、（本题满分 6 分）设 \(\displaystyle f(x)=\begin{cases}\dfrac{\displaystyle g(x)-e^{-x}}{\displaystyle x},\allowbreak &x\ne0,\allowbreak \\0,\allowbreak &x=0,\allowbreak \end{cases}\) 其中 \(\displaystyle g(x)\) 有二阶连续导数，且 \(\displaystyle g(0)=1,\allowbreak g^{\prime}(0)=-1\)。（1）求 \(\displaystyle f^{\prime}(x)\)；（2）讨论 \(\displaystyle f^{\prime}(x)\) 在 \(\displaystyle ( -\infty,\allowbreak +\infty)\) 上的连续性。


四、（本题满分 6 分）设函数 \(\displaystyle z=f(u)\)，方程 \(\displaystyle u=\varphi(u)+\displaystyle \int_y^x p(t)\,dt\) 确定 \(\displaystyle u\) 是 \(\displaystyle x,\allowbreak y\) 的函数，其中 \(\displaystyle f(u),\allowbreak \varphi(u)\) 可微；\(\displaystyle p(t),\allowbreak \varphi^{\prime}(u)\) 连续，且 \(\displaystyle \varphi^{\prime}(u)\ne1\)。求 \(\displaystyle p(y)\dfrac{\displaystyle \partial z}{\displaystyle \partial x}+p(x)\dfrac{\displaystyle \partial z}{\displaystyle \partial y}\)。


五、（本题满分 6 分）计算 \(\displaystyle \int_0^{+\infty}\dfrac{\displaystyle xe^{-x}}{\displaystyle (1+e^{-x})^2}\,dx\)。


六、（本题满分 5 分）设 \(\displaystyle f(x)\) 在区间 \(\displaystyle [0,\allowbreak 1]\) 上可微，且满足条件 \(\displaystyle f(1)=2\displaystyle \int_0^{1/2}xf(x)\,dx\)。试证：存在 \(\displaystyle \xi\in(0,\allowbreak 1)\)，使 \(\displaystyle f(\xi)+\xi f^{\prime}(\xi)=0\)。


七、（本题满分 6 分）设某种商品的单价为 \(\displaystyle p\) 时，售出的商品数量 \(\displaystyle Q\) 可以表示成 \(\displaystyle Q=\dfrac{\displaystyle a}{\displaystyle p+b}-c\)，其中 \(\displaystyle a,\allowbreak b,\allowbreak c\) 均为正常数，且 \(\displaystyle a>bc\)。（1）求 \(\displaystyle p\) 在何范围变化时，使相应销售额增加或减少；（2）要使销售额最大，商品单价 \(\displaystyle p\) 应取何值？最大销售额是多少？


\subsection*{试卷 IV（续）}

八、（本题满分 6 分）求微分方程 \(\displaystyle \dfrac{\displaystyle dy}{\displaystyle dx}=\dfrac{\displaystyle y-\sqrt{x^2+y^2}}{\displaystyle x}\) 的通解。


九、（本题满分 8 分）设矩阵 \(\displaystyle A=\begin{pmatrix}0&1&0&0\\1&0&0&0\\0&0&y&1\\0&0&1&2\end{pmatrix}\)。（1）已知 \(\displaystyle A\) 的一个特征值为 3，试求 \(\displaystyle y\)；（2）求可逆矩阵 \(\displaystyle P\)，使 \(\displaystyle (AP)^{\mathrm T}(AP)\) 为对角矩阵。


十、（本题满分 8 分）设向量组 \(\displaystyle \alpha_1,\allowbreak \alpha_2,\allowbreak \ldots,\allowbreak \alpha_t\) 是齐次线性方程组 \(\displaystyle Ax=0\) 的一个基础解系，向量 \(\displaystyle \beta\) 不是方程组 \(\displaystyle Ax=0\) 的解，即 \(\displaystyle A\beta\ne0\)。试证明：向量组 \(\displaystyle \beta,\allowbreak \beta+\alpha_1,\allowbreak \beta+\alpha_2,\allowbreak \ldots,\allowbreak \beta+\alpha_t\) 线性无关。


十一、（本题满分 7 分）假设一部机器在一天内发生故障的概率为 0.2，机器发生故障时全天停止工作。若一周 5 个工作日里无故障，可获利润 10 万元；发生一次故障仍可获利润 5 万元；发生二次故障所获利润 0 元；发生三次或三次以上故障就要亏损 2 万元。求一周内利润的期望是多少？


十二、（本题满分 6 分）考虑一元二次方程 \(\displaystyle x^2+Bx+C=0\)，其中 \(\displaystyle B,\allowbreak C\) 分别是将一枚色子（骰子）接连掷两次先后出现的点数，求该方程有实根的概率 \(\displaystyle p\) 和有重根的概率 \(\displaystyle q\)。


十三、（本题满分 6 分）假设 \(\displaystyle X_1,\allowbreak X_2,\allowbreak \ldots,\allowbreak X_n\) 是来自总体 \(\displaystyle X\) 的简单随机样本。已知 \(\displaystyle E(X^k)=\alpha_k\)（\(\displaystyle k=1,\allowbreak 2,\allowbreak 3,\allowbreak 4\)）。证明：当 \(\displaystyle n\) 充分大时，随机变量 \(\displaystyle Z_n=\dfrac{\displaystyle 1}{\displaystyle n}\displaystyle \sum_{i=1}^nX_i^2\) 近似服从正态分布，并指出其分布参数。


\subsection*{1996 年全国硕士研究生招生考试试题（试卷 V）}

\subsection*{一、填空题（本题共 5 小题，每小题 3 分，满分 15 分）}

（1）【同试卷 IV 第一、（1）题】


（2）【同试卷 IV 第一、（2）题】


（3）设 \(\displaystyle y=\ln\left(x+\sqrt{1+x^2}\right)\)，则 \(\displaystyle \left.y^{\prime\prime\prime}\right|_{x=\sqrt3}=\)\_\_\_\_\_\_。


（4）5 阶行列式 \(\displaystyle \begin{vmatrix}1-a&a&0&0&0\\-1&1-a&a&0&0\\0&-1&1-a&a&0\\0&0&-1&1-a&a\\0&0&0&-1&1-a\end{vmatrix}=\)\_\_\_\_\_\_。


\subsection*{试卷 V（续）}

一、（5）一实习生用同一台机器接连独立地制造 3 个同种零件，第 \(\displaystyle i\) 个零件是不合格品的概率 \(\displaystyle P_i=\dfrac{\displaystyle 1}{\displaystyle i+1}\)（\(\displaystyle i=1,\allowbreak 2,\allowbreak 3\)），以 \(\displaystyle X\) 表示 3 个零件中合格品的个数，则 \(\displaystyle P\{X=2\}=\)\_\_\_\_\_\_。


\subsection*{二、选择题（本题共 5 小题，每小题 3 分，满分 15 分）}

（1）设 \(\displaystyle f^{\prime}(x_0)=f^{\prime\prime}(x_0)=0,\allowbreak  f^{\prime\prime\prime}(x_0)>0\)，则下列选项正确的是（　）。


（A）\(\displaystyle f^{\prime}(x_0)\) 是 \(\displaystyle f^{\prime}(x)\) 的极大值。 （B）\(\displaystyle f(x_0)\) 是 \(\displaystyle f(x)\) 的极大值。 （C）\(\displaystyle f(x_0)\) 是 \(\displaystyle f(x)\) 的极小值。 （D）\(\displaystyle (x_0,\allowbreak f(x_0))\) 是曲线 \(\displaystyle y=f(x)\) 的拐点。


（2）设 \(\displaystyle f(x)\) 处处可导，则（　）。


（A）当 \(\displaystyle \lim_{x\to+\infty}f^{\prime}(x)=+\infty\) 时，必有 \(\displaystyle \lim_{x\to+\infty}f(x)=+\infty\)。 （B）当 \(\displaystyle \lim_{x\to+\infty}f(x)=+\infty\) 时，必有 \(\displaystyle \lim_{x\to+\infty}f^{\prime}(x)=+\infty\)。


（C）当 \(\displaystyle \lim_{x\to-\infty}f^{\prime}(x)=-\infty\) 时，必有 \(\displaystyle \lim_{x\to-\infty}f(x)=-\infty\)。 （D）当 \(\displaystyle \lim_{x\to-\infty}f(x)=-\infty\) 时，必有 \(\displaystyle \lim_{x\to-\infty}f^{\prime}(x)=-\infty\)。


（3）【同试卷 IV 第二、（3）题】


（4）【同试卷 IV 第二、（4）题】


（5）设 \(\displaystyle A,\allowbreak B\) 为任意两个事件，且 \(\displaystyle A\subset B,\allowbreak  P(B)>0\)，则下列选项必然成立的是（　）。


（A）\(\displaystyle P(A)<P(A\mid B)\)。 （B）\(\displaystyle P(A)\leq P(A\mid B)\)。 （C）\(\displaystyle P(A)>P(A\mid B)\)。 （D）\(\displaystyle P(A)\geq P(A\mid B)\)。


三、（本题满分 6 分）【同试卷 IV 第三题】


四、（本题满分 7 分）设 \(\displaystyle f(x,\allowbreak y)=\displaystyle \int_0^{xy}e^{-t^2}\,dt\)，求 \(\displaystyle \dfrac{\displaystyle x}{\displaystyle y}\dfrac{\displaystyle \partial^2f}{\displaystyle \partial x^2}-2\dfrac{\displaystyle \partial^2f}{\displaystyle \partial x\partial y}+\dfrac{\displaystyle y}{\displaystyle x}\dfrac{\displaystyle \partial^2f}{\displaystyle \partial y^2}\)。


五、（本题满分 6 分）【同试卷 IV 第五题】


六、（本题满分 7 分）【同试卷 IV 第六题】


七、（本题满分 9 分）已知一抛物线通过 \(\displaystyle x\) 轴上的两点 \(\displaystyle A(1,\allowbreak 0),\allowbreak B(3,\allowbreak 0)\)。（1）求证：两坐标轴与该抛物线所围图形的面积等于 \(\displaystyle x\) 轴与该抛物线所围图形的面积；（2）计算上述两平面图形绕 \(\displaystyle x\) 轴旋转一周所产生的两个旋转体体积之比。


八、（本题满分 6 分）设 \(\displaystyle f(x)\) 在 \(\displaystyle [a,\allowbreak b]\) 上连续，在 \(\displaystyle (a,\allowbreak b)\) 内可导，且 \(\displaystyle \dfrac{\displaystyle 1}{\displaystyle b-a}\displaystyle \int_a^b f(x)\,dx=f(b)\)。求证：在 \(\displaystyle (a,\allowbreak b)\) 内至少存在一点 \(\displaystyle \xi\)，使 \(\displaystyle f^{\prime}(\xi)=0\)。


\subsection*{试卷 V（续）}

九、（本题满分 9 分）已知线性方程组 \(\displaystyle \begin{cases}x_1+x_2-2x_3+3x_4=0,\allowbreak \\2x_1+x_2-6x_3+4x_4=-1,\allowbreak \\3x_1+2x_2+px_3+7x_4=-1,\allowbreak \\x_1-x_2-6x_3-x_4=t.\end{cases}\) 讨论参数 \(\displaystyle p,\allowbreak t\) 取何值时，方程组有解、无解；当有解时，试用其导出组的基础解系表示其通解。


十、（本题满分 7 分）设有 4 阶方阵 \(\displaystyle A\) 满足条件 \(\displaystyle |3E+A|=0,\allowbreak \ AA^{\mathrm T}=2E,\allowbreak \ |A|<0\)，其中 \(\displaystyle E\) 是 4 阶单位阵。求方阵 \(\displaystyle A\) 的伴随矩阵 \(\displaystyle A^*\) 的一个特征值。


十一、（本题满分 7 分）【同试卷 IV 第十一题】


十二、（本题满分 6 分）假设一电路装有三个同种电气元件，其工作状态相互独立，且无故障工作的时间都服从参数为 \(\displaystyle \lambda>0\) 的指数分布。当三个元件都无故障时，电路正常工作，否则整个电路不能正常工作。试求电路正常工作的时间 \(\displaystyle T\) 的概率分布。

\end{document}
